\section{Experimental design}
\label{sec:experiments}

The evaluation has three parallel parts: core results, ablation studies, and
analysis. Domains enter confirmatory experiments only after a readiness audit.
The complete domain panorama is retained independently of the selected execution
scope. Models, datasets, human protocols, budgets, and primary contrasts are
frozen before final outcomes are observed.

\subsection{Core results}
\label{sec:core}

\begin{figure}[t]
\centering
\begin{tikzpicture}[line cap=round,line join=round]
  \path[use as bounding box] (0,0) rectangle (13.7,5.78);
  \foreach \x/\t/\col in {0/{(a) C1: alignment and selection}/pBlue,
    4.65/{(b) C2: critic-guided refinement}/pTeal,
    9.30/{(c) C3: reward training (planned)}/pPlum}{
    \node[pTitle,text=\col,anchor=west] at (\x,5.55) {\t};
    \fill[\col!4,rounded corners=3pt] (\x,.75) rectangle (\x+4.4,5.26);
  }

  \node[pText,align=center] at (2.2,4.97) {Same input and fixed pool};
  \foreach \x/\n in {.8/1,1.73/2,2.66/3,3.59/4}{
    \draw[pBlue!60,fill=white,rounded corners=1pt] (\x-.28,4.13) rectangle (\x+.28,4.64);
    \node[pSmall,text=pBlue] at (\x,4.39) {$y_{\n}$};
  }
  \node[pCard,draw=pBlue!40,text width=3.68cm,minimum height=.55cm,align=center]
    (judges) at (2.2,3.56) {Frozen evaluators\\Identical scoring protocol};
  \draw[pArrow] (2.2,4.08) -- (judges.north);
  \node[pCard,draw=pBlue!40,text width=1.36cm,align=center,minimum height=.62cm]
    (selects) at (1.06,2.40) {Baseline\\selection};
  \node[pCard,draw=pBlue,text width=1.36cm,align=center,minimum height=.62cm]
    (selectp) at (3.34,2.40) {P2E\\selection};
  \draw[pSoftArrow] (judges.south) -- (2.2,3.04) -| (selects.north);
  \draw[pSoftArrow] (2.2,3.04) -| (selectp.north);
  \node[pSmall,align=center] at (2.2,1.30)
    {Full human rankings\\Agreement / selection / regret};
  \draw[pSoftArrow] (selects.south) -- (1.06,1.76);
  \draw[pSoftArrow] (selectp.south) -- (3.34,1.76);

  \node[pCard,draw=pTeal,fill=white,text width=2.85cm,minimum height=.52cm,align=center]
    (y0) at (6.85,4.64) {Shared fresh $y_0$};
  \foreach \x/\id/\t in {5.67/cs/Seed critic,8.03/cp/P2E critic}{
    \node[pCard,draw=pTeal!50,text width=1.57cm,minimum height=.54cm,align=center]
      (\id) at (\x,3.61) {\t};
    \draw[pSoftArrow] (y0.south) -- (6.85,4.10) -| (\id.north);
  }
  \foreach \x/\id/\t in {5.67/ys/$y_{\mathrm{Seed}}'$,8.03/yp/$y_{\mathrm{P2E}}'$}{
    \node[pCard,draw=pTeal!55,fill=pTeal!6,minimum width=1.22cm,minimum height=.55cm,align=center]
      (\id) at (\x,2.28) {\t};
  }
  \draw[pSoftArrow] (cs.south) -- (ys.north);
  \draw[pSoftArrow] (cp.south) -- (yp.north);
  \node[pSmall,align=center,fill=white,inner sep=2pt] at (6.85,2.96)
    {Same reviser and budget};
  \node[pSmall,align=center] at (6.85,1.28)
    {New blind ranking: $y_0$, $y_{\mathrm{Seed}}'$, $y_{\mathrm{P2E}}'$\\
     Primary: P2E vs.\ Seed};

  \node[pText,align=center] at (11.50,4.95) {Same base and inputs; frozen rewards};
  \node[pSmall,align=center] at (11.50,4.67) {Base arm: no parameter update};
  \foreach \x/\id/\t in {10.40/rs/Seed reward,12.60/rp/P2E reward}{
    \node[pCard,draw=pPlum!50,text width=1.60cm,minimum height=.61cm,align=center]
      (\id) at (\x,4.15) {\t};
  }
  \foreach \x/\id in {10.40/ps,12.60/pp}{
    \node[pCard,draw=pPlum!55,fill=pPlum!6,text width=1.45cm,minimum height=.68cm,align=center]
      (\id) at (\x,2.98) {GRPO\\(LoRA)};
    \draw[pSoftArrow] (\x,3.81) -- (\id.north);
    \draw[pSoftArrow] (\id.south) -- (\x,2.25);
    \draw[pSoftArrow] (\x,2.25) -- (\x+.94,2.25) -- (\x+.94,3.0) -- (\id.east);
    \node[pSmall,align=center] at (\x,1.91) {Own rollouts};
  }
  \node[pSmall,align=center] at (11.50,1.14)
    {Base, Seed-GRPO, P2E-GRPO\\
     New outputs and blinded $H_0/H_1$};

  \draw[pLine,line width=.6pt] (0,.52) -- (13.7,.52);
  \node[pSmall,align=center] at (6.85,.23)
    {Independent ID / OOD assessment \quad $\vert$ \quad External Metric Reserve \quad $\vert$ \quad Original-input grouped uncertainty};
\end{tikzpicture}

\caption{Three controlled uses of the evolved evaluator. Panel (a), C1, holds the
candidate pool fixed and compares ranking and selection. Panel (b), C2, branches
from a freshly generated shared $y_0$ and changes only the critic. Panel (c), C3,
contrasts an unchanged Base policy with two policies trained under frozen rewards.
The trained policies use matched inputs and compute but generate their own rollouts.
Final judgments are blinded
and external metrics are reserved from evolution. C3 is a planned training study,
not an executed experiment.}
\label{fig:downstream}
\end{figure}

\paragraph{C1: evaluator alignment and selection.}
Frozen P2E evaluators (Figure~\ref{fig:downstream}a) are compared with the best single metric, a metric ensemble,
a strong static judge, and a static tool-augmented judge. Prompt-optimization
and program-search controls use declared versions and comparable edit interfaces.
Matched comparisons share the base model, human supervision, inputs, candidate
pools, and scoring protocol. Search cost and deployment cost are reported
separately.

For each test input, all evaluators receive the same candidate group, with
a default design of $N=4$. Complete human rankings support global pairwise
agreement, top-region agreement involving the human top-two candidates, Best-of-$N$
selection, and selection regret. When only ranks are available, normalized regret
is $(\operatorname{rank}(\hat y)-1)/(N-1)$ for the selected output $\hat y$.
A pair-only dataset supports pairwise evaluation, not an $N$-way selection claim.
Core $H_1$, worst-stratum performance, coverage, variance, and cost accompany
overall agreement.

ID and OOD results are reported separately. OOD changes dataset or source
distribution while preserving the evaluation standard. A difficult subset alone
is a difficulty test, not an OOD source. Final sampling is independent of mining
conflicts. All splitting and statistical resampling preserve original input/source
groups. Paired effects use grouped 95\% intervals. Existing annotations are reused
only when their construct and candidate groups match the endpoint.

\paragraph{C2: critic-guided refinement.}
For each unseen input, a frozen generator produces one fresh output $y_0$.
The seed critic and P2E critic each inspect that same output and permitted evidence,
then return the same critique format. A common generator performs one refinement step
under identical prompts, decoding, tool access, and token budgets
(Figure~\ref{fig:downstream}b). A separate refinement-validation set determines
these settings before ID/OOD generation.
New blinded human rankings assess $y_0$, Seed-refined, and P2E-refined.
The primary contrast is P2E-refined versus Seed-refined, accompanied by improvement
and regression rates against $y_0$, core dimensional judgments, and reserved
external metrics. Labels on pre-existing outputs cannot assess these new outputs.

\paragraph{C3: reward-guided training.}
The planned training study compares Base, Seed-GRPO, and P2E-GRPO from the same
initial checkpoint (Figure~\ref{fig:downstream}c), using LoRA and group relative policy optimization
\citep{hu2022lora,shao2024deepseekmath}. The two training arms share prompts and
their order, optimizer settings, group size, token limits, update count, and
policy-training compute. Rewards are frozen before training. Rollouts are sampled
online by each arm's own evolving policy and are not shared. Reward-call counts
and token limits are matched; actual scoring cost is reported separately.
Frozen final checkpoints generate new ID/OOD outputs for new blinded $H_0/H_1$
and reserved external metrics. Checkpoint selection, invalid-output policy, and
reward-hacking diagnostics are specified in Appendix~\ref{app:training}.
C3 remains a design-only study here; no parameter training or training benefit
is reported.

\subsection{Ablation studies}
\label{sec:ablations}

The ablations separate annotation choices from changes to the evaluator and search
procedure (Table~\ref{tab:ablations}). Mining strategies are compared at equal
human budgets $N_H\in\{20,50,100,200\}$, measured in original-input annotation
units. The endpoint is final held-out alignment after evolution, not the fraction
of errors found during acquisition. When feedback granularity changes annotation
time, both unit count and time/cost are reported. Component controls preserve
unaffected interfaces; an invalid dependency is reported instead of silently
reverting the entire evaluator to its seed.

\begin{table}[t]
\caption{Controlled ablation families. Each row tests a different hypothesis;
none is a completed result. The full condition matrix is in
Appendix~\ref{app:ablations}.}
\label{tab:ablations}
\centering\small
\begin{tabularx}{\linewidth}{@{}P{2.3cm}YY@{}}
\toprule
Variable & Principal conditions & Fixed quantities \\
\midrule
Preference mining & Random; uncertainty; model conflict; instability; P2E &
Human budget, evolution procedure, test set \\
Feedback detail & $H_0$; $+H_1$; $+H_2$; $+H_3$ &
Same pairs and annotation provenance \\
Editable components & Full; rubric/procedure only; fixed tools/routes; routing only &
Seed, model, opportunity and measurement budgets \\
Search allocation & Independent proposals; history only; direction continuation; pause/resume &
Initial state, evidence access, total compute \\
Component attribution & Tool on/off; evolved/seed tool; adaptive/static route; fixed fusion &
Other components and paired input groups \\
\bottomrule
\end{tabularx}
\end{table}

\subsection{Analysis and studies}
\label{sec:analysis}

Trajectory analysis reports accepted directions, implementation attempts,
resource use, and the evidence supporting each scoped refutation.
Failure-regime analysis compares rubric, tool, routing, and fusion contributions
with the remaining gap to human agreement. A domain is described as
specification-limited, observation-limited, or headroom-limited only after these
comparisons; strata within one domain may differ.
Feedback-efficiency analysis relates held-out changes to annotation units, time,
and granularity. Adapter analysis compares direct human judgments with labels
derived from rankings, pointwise scores, or error annotations on the same subset.
Repeated scoring and independent search runs are separate sources of variation.
